\chapter{Related Work}
\label{ch:relatedWork}

\section{Argument Mining on Deliberation Platforms}
\label{sec:rw-am-deliberation}

\textcite{elguendouze2025am4dsp} present AM4DSP, a system that performs argumentation mining
directly on structured decentralized discussion platforms built for deliberative democracy,
classifying statements into position and evidence roles and identifying relations between them,
evaluated on the BCause platform. This is close precedent for the relation-extraction stage of
this thesis: both projects mine argumentative structure from real or realistic civic discussion
data rather than curated debate corpora. The key divergence is scope. AM4DSP targets
\emph{discussion-level} analysis within a single deliberation thread, structuring the argumentation
inside one conversation; this thesis instead asks a \emph{retrieval} question across an entire
corpus of thousands of independent arguments spanning many policies, where the challenge is
selecting and organizing a relevant, non-redundant, conflicting subset rather than parsing the
structure of one discussion. AM4DSP's mined relations are consumed by human moderators for
sense-making; the relations extracted in this thesis are consumed by an automated retrieval and
summarization pipeline evaluated against graded relevance judgments, a different downstream use
that motivates a different evaluation protocol entirely (Chapter~\ref{ch:evaluation}).

\textcite{anastasiou2025bcause} describe BCause, a human-AI collaboration tool that supports hybrid
mapping and ideation in argumentation-grounded deliberation, funded under the EU ORBIS programme
on augmenting trust and transparency in deliberative democracy. Like this thesis, BCause is
built on the premise that making argumentative structure explicit improves a deliberation
process; unlike this thesis, BCause is fundamentally a human-in-the-loop authoring and mapping
tool, where a facilitator or participant actively builds and edits the argument map. This thesis
instead evaluates a fully automated retrieval pipeline with no human curation step at query time,
trading the higher precision a human curator could provide for the ability to operate at the
scale of a 30{,}500-argument corpus, where manual mapping is not feasible.

\textcite{behrendt2025survey} survey the broader landscape of machine-learning methods for
enhancing deliberation in online political discussions, cataloguing which tasks current AI
methods address and identifying open challenges. Their survey confirms that argument-relation
detection and redundancy/conflict identification, the two relation types this thesis constructs,
are recognized open problems in the deliberation-technology literature rather than a solved
prerequisite; this motivates treating the calibration of these relations (Chapter~\ref{ch:methodology})
as a contribution in its own right, not a routine preprocessing step.

\section{Knowledge Graphs for Argumentative and Deliberative Data}
\label{sec:rw-kg-deliberation}

\textcite{plenz2024pakt} construct PAKT, a perspectivized argumentation knowledge graph built
directly for deliberation analysis, encoding not only argumentative relations but the
perspective, or stance, from which each argument is made. PAKT's graph construction is close in
spirit to this thesis's own \textsc{contradicts}/\textsc{equivalent} edges, and its explicit
handling of perspective anticipates a challenge this thesis meets independently on the Grand
D\'ebat National corpus, where stance and topic are shown to be entangled in embedding space
(Section~\ref{sec:gdn-stance}). The divergence is again one of downstream purpose: PAKT is
positioned as an analysis and visualization tool for researchers studying deliberation, not as
the retrieval backend for a conflict-mapping or summarization system, and its evaluation is
correspondingly framed around graph quality and analytical utility rather than retrieval metrics
such as nDCG or diversity.

\textcite{vagnoni2025delibkg} introduce the Deliberation Knowledge Graph, integrating deliberation
data from European Parliament proceedings, civic participation platforms, and public forums under
a shared ontology, and demonstrate its use for examining argument quality and tracing the
evolution of policy positions across institutional and civic spheres. This work shares this
thesis's premise that an explicit graph over argumentative data supports analysis that flat text
cannot, but it is oriented toward \emph{integration} across heterogeneous, multi-source
institutional data under a common ontology, whereas this thesis constructs a single-source,
within-corpus graph and asks a comparative retrieval question (graph-aware versus
embedding-only retrieval) that the Deliberation Knowledge Graph's descriptive, ontology-integration
framing does not address.

\textcite{debatekg2023} build argumentative semantic knowledge graphs over a large policy-debate
corpus and construct debate cases via constrained shortest-path traversal over the resulting
graph, evaluated on an extension of the DebateSum dataset, a close relative of the IBM Project
Debater corpus used in this thesis. DebateKG is the closest existing work to this thesis in terms
of \emph{data domain}: both operate on competitive or crowd-authored policy-debate text at
comparable scale. The divergence is in task and evaluation. DebateKG's shortest-path traversal
is optimized to construct a single coherent debate case (a chain of connected arguments serving
one side), not to retrieve a balanced, graded-relevant, non-redundant set spanning both PRO and
CON stances for a conflict map, and DebateKG reports no comparison against an embedding-only
baseline, leaving open exactly the empirical question, whether graph traversal outperforms dense
retrieval on this class of corpus, that this thesis investigates directly (RQ1,
Section~\ref{sec:rq}).

\section{RAG and Graph-RAG in Civic and Governmental Applications}
\label{sec:rw-graphrag-civic}

\textcite{papageorgiou2025egovgraphrag} propose a multi-agent GraphRAG framework for e-government,
motivated explicitly by the observation that flat, unstructured RAG retrieval limits multi-hop
reasoning and interpretability on structured government data, and demonstrate the framework on a
policy-focused question-answering case study. This is the closest existing work to this thesis's
core empirical question, a direct application of the RAG-versus-Graph-RAG comparison to policy
text. It differs from this thesis in three respects that matter for what can be concluded from
it. First, its graph is built over policy \emph{documents} (structured e-government datasets,
press releases), not over individual \emph{citizen arguments}, so its edges represent document or
entity relations rather than the contradiction and equivalence relations between argumentative
claims that this thesis's retrieval systems exploit. Second, its evaluation is a single case
study demonstrating the framework's operation, not a controlled comparison against a matched
embedding-only baseline under a shared evaluation protocol, so no quantitative claim is made
about whether the graph-aware component actually outperforms flat retrieval on this data, the
precise gap this thesis's UC1 evaluation is designed to close. Third, it does not evaluate a
diversity or redundancy criterion, whereas avoiding redundant retrieval is central to this
thesis's conflict-map task (UC1). Taken together, this thesis can be read as supplying the
controlled, quantitative retrieval evaluation that this line of applied GraphRAG work motivates
but does not itself provide.

\textcite{tessler2024science} present the Habermas Machine, an LLM-based system that generates
group statements intended to represent common ground across a set of citizen opinions, evaluated
in large-scale deliberation experiments. While not a retrieval system and not graph-based, this
work establishes that LLM-mediated systems can measurably improve deliberative outcomes at scale,
which motivates this thesis's underlying premise that the retrieval stage feeding such a
generation system is worth optimizing carefully; a summarization or common-ground system is only
as good as the arguments it is given to work with, a claim this thesis tests directly in UC2,
where retrieval-quality differences are evaluated for whether they survive the summarization
bottleneck (Section~\ref{sec:uc2-methodology}).